\section{IBM Quantum Hardware Experimental Validation (Implementation)}
Hardware experiments were chosen only where physical execution answers a question that simulation cannot: how far real-device behaviour departs from ideal and backend-derived noisy simulation for the project's quantum components. The single-qubit BB84 preparation and measurement circuits and the intercept--resend circuit model are small enough to run on current hardware, and they connect directly to Layer~2; a two-qubit Grover circuit shows the effect of two-qubit gates. The backend was \HwBackend{} (pilot job \HwPilotJob{}, then \HwJobs{} independent jobs of 52 circuits at \HwShots{} shots). The scripts \texttt{scripts/ibm\_quantum/run\_hardware\_campaign.py} (submission), \texttt{run\_aer\_controls.py} (ideal and noisy Aer controls with the same transpiled circuits) and \texttt{analyze\_campaign.py} (analysis from saved files only) implement the study; credentials are read from a gitignored environment file and never stored. Qubits were selected from the calibration snapshot by a fixed rule before any job ran, and the snapshot, raw counts, job identifiers and transpilation statistics are saved under \texttt{results/ibm\_quantum/}. Results are given in the chapter \emph{IBM Quantum Hardware Experimental Validation}.
